\section{Obstruction Theory \& Characteristic Classes}
\label{sec:ObstructionTheory&CharacteristicClasses}
\subsection{A Brief Introduction to Obstruction Theory}
In homotopy theory, one often encounters the following problem: 
\begin{center}
\fbox{
    \parbox{0.9\linewidth}{
        Given a pair of spaces $(X,A)$ and a map $f:A\to Y$. Can $f$ be extended 
        to a map $F:X\to Y$?
    }
}
\end{center}
Obstruction theory provides a systematic way to study 
this problem. 

$\mathbf{Toy\ Model:}$ A map $f:S^n\to Y$ can be extended to a map 
$F:D^{n+1}\to Y$ if and only if $f$ is null-homotopic.

Now we consider the general case. Let $K$ be a CW complex, $Y$ be a simple 
space, i.e., $\pi_1(Y)$ acts trivially on $\pi_n(Y)$ for all $n$. Under this 
condition, we can unambiguously use the notation $\pi_n(Y) = [S^n, Y]$.  
Suppose we have a map $f:K^{(n)}\to Y$, where $K^{(n)}$ is the 
$n$-skeleton of $K$. We want to extend $f$ to a map $F:K^{(n+1)}\to Y$. 
To do this, we only need to extend $f$ to each $(n+1)$-cell $e_\alpha^{n+1}$. 

Suppose the attaching map of $e_\alpha^n$ is 
$\varphi_\alpha:S^{n}\to K^{(n)}.$ Then $f\circ \varphi_\alpha$ is a map
from $S^{n}$ to $Y$. By the toy model, $f\circ \varphi_\alpha$ can be
extended to $D^{n+1}$ if and only if $[f\circ \varphi_\alpha] = 0$ in
$\pi_{n}(Y)$. We in fact get a cochain 
$c^{n+1}(f)\in C^{n+1}_{cell}(K;\pi_{n}(Y))$: 
\begin{align*}
    c^{n+1}(f):C_{n+1}^{cell}(K)\to \pi_{n}(Y) \\
    e_\alpha^{n+1}\mapsto [f\circ \varphi_\alpha].
\end{align*}
We call $c^{n+1}(f)$ the \textbf{${(n+1)}$-th obstruction cochain} of $f$. 
There are some basic facts about $c^{n+1}(f)$: 
\begin{itemize}
    \item $f$ can be extended to $K^{(n+1)}$ if and only if $c^{n+1}(f) = 0$; 
    \item If $f,f':K^{(n)}\to Y$ are homotopic, then 
    $c^{n+1}(f) = c^{n+1}(f')$; 
    \item (Natruality) Suppose $L$ is another CW complex and $\phi:L\to K$ 
    is a cellular map, then $c^{n+1}(f\circ \phi) = \phi^{\sharp}c^{n+1}(f)$.
\end{itemize}

At the following, we give some important properties of $c^{n+1}(f)$ without 
proof. 
\begin{theorem}
    $c^{n+1}(f)$ is a cocycle, i.e., $\delta c^{n+1}(f) = 0$. Thus it 
    defines a cohomology class $[c^{n+1}(f)]\in H^{n+1}(K;\pi_{n}(Y))$. 
    We call $[c^{n+1}(f)]$ the \textbf{$(n+1)$-th obstruction class} of $f$.
\end{theorem}

A similar concept is defined for homotopies. Suppose we have two maps
$f_0,f_1:K^{(n)}\to Y$ and a homotopy $h:K^{(n-1)}\times I\to Y$ between
$f_0|_{K^{(n-1)}}$ and $f_1|_{K^{(n-1)}}$. We want to extend $h$ to a 
homotopy $H:K^{(n)}\times I\to Y$ between $f_0$ and $f_1$. To do this, 
we only need to extend $h$ to each $n$-cell $e_\alpha^n$. Let 
$\varphi_\alpha:S^{n-1}\to K^{(n-1)}$ be the attaching map of $e_\alpha^n$, 
$\psi_\alpha:D^n\to K^{(n)}$ be the characteristic map of $e_\alpha^n$. Then 
we can form a map
\begin{align*}
    h_\alpha:(D^n\times \{0,&1\})\cup(\p D^n\times I)\to Y \\ 
    (x,0)&\mapsto f_0\circ \psi_\alpha(x) \\
    (x,1)&\mapsto f_1\circ \psi_\alpha(x) \\
    (x,t)&\mapsto h\circ (\varphi_\alpha(x),t).
\end{align*}
Since $(D^n\times \{0,1\})\cup(\p D^n\times I)\cong S^n$, $h_\alpha$ defines
a map from $S^n$ to $Y$. It can extend to $D^n\times I\cong D^{n+1}$ if and 
only if $[h_\alpha] = 0$ in $\pi_n(Y)$. We in fact get a cochain
$c^n(f_0,f_1,h)\in C^n_{cell}(K;\pi_n(Y))$:
\begin{align*}
    c^n(f_0,f_1,h):C_n^{cell}(K) & \to \pi_n(Y) \\
    e_\alpha^n & \mapsto [h_\alpha]. 
\end{align*} 
We call $c^n(f_0,f_1,h)$ the \textbf{$n$-th deformation cochain} of 
$(f_0,f_1,h)$. If $f_0|_{K^{(n-1)}} = f_1|_{K^{(n-1)}}$ and $h$ is the 
trivial homotopy, we simply write $c^n(f_0,f_1)$ for $c^n(f_0,f_1,h)$ and 
call it the \textbf{$n$-th difference cochain} of $(f_0,f_1)$. 
There are some basic facts about $c^n(f_0,f_1,h)$:
\begin{itemize}
    \item $h$ can be extended to a homotopy $H:K^{(n)}\times I\to Y$ between 
    $f_0$ and $f_1$ if and only if $c^n(f_0,f_1,h) = 0$; 
    % \item If $h,h':K^{(n-1)}\times I\to Y$ are homotopic rel. boundary, then 
    % $c^n(f_0,f_1,h) = c^n(f_0,f_1,h')$; 
    \item (Natruality) Suppose $L$ is another CW complex and $\phi:L\to K$ 
    is a cellular map, then 
    $c^n(f_0\circ \phi,f_1\circ \phi,h\circ(\phi\times \id_I)) = 
    \phi^{\sharp}c^n(f_0,f_1,h)$.
    \item $c^n(f_0,f_2,h*h') = c^n(f_0,f_1,h) + c^n(f_1,f_2,h')$, where
    $h*h'$ is the concatenation of $h$ and $h'$. Especially, if 
    $f_0|_{K^{(n-1)}} = f_1|_{K^{(n-1)}} = f_2|_{K^{(n-1)}}$ and $h,h'$ are
    trivial homotopies, then $c^n(f_0,f_2) = c^n(f_0,f_1) + c^n(f_1,f_2)$.
\end{itemize}

In general, $c^n(f_0,f_1,h)$ is not a cocycle. But we have the following
formula. 
\begin{theorem}
    $\delta c^n(f_0,f_1,h) = c^{n+1}(f_1) - c^{n+1}(f_0)$.
\end{theorem} 

There is a useful lemma about difference cochains.
\begin{lemma}
    For any $f:K^{(n)}\to Y$ and any cochain $d\in C^n_{cell}(K;\pi_n(Y))$, 
    there exists a map $f':K^{(n)}\to Y$ such that 
    $f'|_{K^{(n-1)}} = f|_{K^{(n-1)}}$ and $d^n(f,f') = d$.
\end{lemma}

With the above preparations, we have the following theorem:
\begin{theorem}[Eilenberg]
    Given a map $f:K^{(n)}\to Y$, $[c^n(f)] = 0$ in 
    $H^{n+1}(K;\pi_n(Y))$ if and only if $f$ can extend to $K^{(n+1)}$ 
    after having been changed on $K^{(n)}\big\backslash K^{(n-1)}$.
\end{theorem}
\begin{proof}
    $(\Leftarrow):$ If $\exists\,f':K^{(n+1)}\to Y$ such that 
    $f'|_{K^{(n-1)}} = f|_{K^{(n-1)}}$, then 
    \begin{equation*}
        c^{n+1}(f) = c^{n+1}(f) - c^{n+1}(f') = \delta d^n(f',f)
    \end{equation*}
    Thus $c^{n+1}(f)$ is a coboundary. 

    $(\Rightarrow):$ If  $[c^{n+1}(f)]=0$, then $\exists\,d\in C^n(K;\pi_n(Y))$ 
    such that $c^{n+1}(f) = \delta d$. By lemma, there exists 
    $f':K^{(n)}\to Y$ such that $f|_{K^{(n-1)}} = f'|_{K^{(n-1)}}$ and 
    $d^n(f,f') = -d$. So 
    \begin{equation*}
        c^{n+1}(f') = c^{n+1}(f) + \delta d^n(f,f') = \delta d - \delta d = 0.
    \end{equation*}
    Thus $f'$ can be extended to $K^{(n+1)}$.
\end{proof}

\begin{example}
    If the target space $Y$ is $n$-connected, i.e., 
    $\pi_k(Y)=0,\; k=0,1,\dots,n$, then every map $f:K^{(0)}\to Y$ can be 
    extended to $K^{(n+1)}$. Besides, any two maps $f,g:K^{(n)}\to Y$ are 
    homotopic. Indeed, $f|_{K^{(0)}}\simeq g|_{K^{(0)}}$ for the reason that 
    $K$ is path-connected. We denote the homotopy by $h_0$. Then $h_0$ can 
    be extened to $h_n:K^{(n)}\times I\to Y$ since all the deformation cochain 
    vanish.
\end{example}

\begin{corollary}
    If $Y$ is $n$-connected and $\dim K\leq n$, then all the maps from $K$ 
    to $Y$ are null-homotopic.
\end{corollary}

\begin{theorem}[Hopf-Whitehead]
    For any CW complex $K$ with $\dim K\leq n$ and $(n-1)$-connected space $Y$, 
    there is a one-to-one correspondence:
    \begin{equation*}
        [K,Y]\longleftrightarrow H^n(K;\pi_n(Y)).
    \end{equation*}
\end{theorem}
If we weaken the dimensional condition and instead strengthen the target space 
to be an Eilenberg-Maclane space $K(\pi,n)$. Then we will get another theorem: 
\begin{theorem}
    For any CW complex $X$, abelian group $\pi$ and integer $n$, there is a 
    one-to-one correspondence: 
    \begin{equation*}
        [X,K(\pi,n)]\longleftrightarrow H^n(X;\pi).
    \end{equation*}
\end{theorem}
\begin{remark}
    In fact, there is a special class $\alpha$ in 
    \begin{equation*}
        H^n(K(\pi,n);\pi)\cong \Hom(H_n(K(\pi,n)),\pi)\cong \Hom(\pi,\pi)
    \end{equation*}
    which corresponds to the identity map in $\Hom(\pi,\pi)$. The correspondence 
    can be explicitly written by 
    \begin{align*}
        [X,K(\pi,n)]&\to H^n(X;\pi) \\
        f&\mapsto f^*\alpha.
    \end{align*}
\end{remark}

\subsection{Characteristic Classes via Obsruction Theory}
\subsubsection{Cohomology with Local Coefficients}
To discrib characteristic classes as obstruction classes, we need cohomology 
with local coefficients. 
% \begin{center}
%     \fbox{
%         \parbox{0.95\linewidth}{
%             Suppose $X$ is a path-connected space. A \textbf{local coefficient system} 
%             on $X$ is a functor
%             \begin{equation*}
%                 \mathcal{G}:\pi(X)\to \text{(Abelian groups)},
%             \end{equation*}
%             where $\pi(X)$ is the fundamental groupoid of $X$. If we fix a base point
%             $x_0\in X$, then $\mathcal{G}$ is determined by the group 
%             $G = \mathcal{G}(x_0)$ and a representation $\rho:\pi_1(X,x_0)\to \Aut(G)$. 
%             We call $G$ the \textbf{typical group} of $\mathcal{G}$.
%         }
%     }
% \end{center}

\noindent\textbf{Local Coefficient System:}

Suppose $X$ is a path-connected space. For any $x\in X$, let $G_x$ be an 
abelian group attaching to $x$. For any path $\gamma_{xy}$ from $x$ to $y$, 
we have an isomorphism $(\gamma_{xy})_*:G_x\to G_y$. If $(\gamma_{xy})_*$ 
only depends on the homotopy class of $\gamma$ rel. endpoints and satisfies 
the following conditions: 
\begin{equation*}
    (\text{constant path})_* = \Id, \quad
    (\gamma_{xy}*\gamma_{yz})_* = (\gamma_{yz})_*\circ(\gamma_{xy})_*.
\end{equation*}
Then we call $\mathcal{G} = \{G_x\}_{x\in X}$ a 
\textbf{local coefficient system} on $X$. If we fix a base point $x_0\in X$, 
then $\mathcal{G}$ is determined by the group $G = G_{x_0}$ and an action 
$\rho:\pi_1(X,x_0)\to \Aut(G)$. We call $G$ the \textbf{typical group} of 
$\mathcal{G}$. 

\begin{example}
    If $\mathcal{G}$ is a local coefficient system on $X$ with typical group 
    $G$ and trivial action. Then it is the ordinary constant coefficient we 
    are familiar with.
\end{example}

\begin{example}
    If $G_x = \mZ_2$ for all $x\in X$, then since $\Aut(\mZ_2)=\{\Id\}$, 
    the local coefficient system with typical group $\mZ_2$ must be trivial. 
\end{example}

\begin{example}
    If $G_x = \mZ$ for all $x\in X$ and $\pi_1(X,x_0)$ acts on $\mZ$ by
    multiplication of $\pm 1$ according to whether the orientation is
    preserved or reversed along the loop. Then $\mathcal{G}$ is called the 
    \textbf{orientation system} on $X$ and denoted by $\mathrm{Or}$. 
\end{example}

\begin{example}
    Let $p:\tilde{X}\to X$ be the universal covering of $X$. Then 
    $\mathcal{G} = \{H_n(p^{-1}(x))\}_{x\in X}$ is a local coefficient system 
    on $X$ with typical group $H_n(\tilde{X})$ and action induced by the 
    deck transformation group $\pi_1(X,x_0)$. 
\end{example}

\begin{example}
    Let $E\to X$ be a vector bundle with fiber $F$. Then 
    $\mathcal{G} = \{H_n(p^{-1}(x))\}_{x\in X}$ is a local coefficient system 
    on $X$ with typical group $H_n(F)$ and action induced by the structure 
    group of $E$.
\end{example}

\begin{example}
    $\mathcal{G} = \{\pi_n(X,x)\}_{x\in X}$ is a local coefficient system 
    on $X$ with typical group $\pi_n(X,x_0)$ and action induced by the 
    fundamental group $\pi_1(X,x_0)$.
\end{example}

\noindent\textbf{Cohomology with Local Coefficients:}

Let $X$ be a CW complex and $\mathcal{G} = \{G_x\}_{x\in X}$ be a local
coefficient system on $X$. We define the cellular cochain groups with local
coefficients by
\begin{equation*}
    C^n_{cell}(X;\mathcal{G}) = \left\{f:\big\{\text{$n$-cells of $X$}\big\}
    \to \bigsqcup_{x\in X}G_x \,\Big|\, f(e_\alpha^n)\in 
    G_{x_\alpha}\right\}
\end{equation*}
where $x_\alpha$ is a chosen point in the $n$-cell $e_\alpha^n$. The 
coboundary operator 
$\delta:C^n_{cell}(X;\mathcal{G})\to C^{n+1}_{cell}(X;\mathcal{G})$
is defined by 
\begin{equation*}
    (\delta f)(e_\beta^{n+1}) = \sum_\alpha d_{\beta\alpha}
    (\gamma_{\alpha\beta})_*f(e_\alpha^n),
\end{equation*}
where $d_{\beta\alpha}$ is the degree of the map 
$\p e_\beta^{n+1}\to X^{(n)}/(X^{(n)}-e^n_\alpha)\cong S^n$ and 
$\gamma_{\alpha\beta}$ is a path from $x_\alpha$ to $x_\beta$. It can be
verified that $\delta^2=0$. Thus we can define the cohomology groups with 
local coefficients by 
\begin{equation*}
    H^n(X;\mathcal{G}) = \frac{\ker(\delta:C^n_{cell}(X;\mathcal{G})\to 
    C^{n+1}_{cell}(X;\mathcal{G}))}{\mathrm{im}(\delta:C^{n-1}_{cell}(X;
    \mathcal{G})\to C^n_{cell}(X;\mathcal{G}))}.
\end{equation*}

\subsubsection*{The Homotopy Groups of $V_k(\mR^n)$}
\textbf{Fact:} The Stiefel manifold $V_k(\mR^n)$ is $(n-k-1)$-connected and 
the first nontrivial homotopy group is 
\begin{equation*}
    \pi_{n-k}(V_k(\mR^n)) = 
    \begin{cases}
        \mZ,& n-k \text{ is even or } k=1; \\
        \mZ_2,& n-k \text{ is odd and } k>1.
    \end{cases}
\end{equation*}
\begin{remark}
    It is evident that $V_k(\mR^n)$ is a simple space.
\end{remark}

\subsubsection*{Characteristic Classes}
Now let's back to vector bundles. Let $\xi$ be an $\mR^n$-bundle over a CW 
complex $B$. We want to investigate how nontrivial $\xi$ is. To do this, we 
try to find $k$-linearly independent sections of $\xi$. 

% First suppose there 
% is a Euclidean metric $g$ on $\xi$. Then we can talk about orthonormal 
% sections. 
% \begin{remark}
%     Any real vector bundle over a paracompact space admits a 
%     Euclidean metric. Each two Euclidean metrics $g_0, g_1$ are 
%     homotopic through $g_t = (1-t)g_0 + tg_1$. 
% \end{remark}

Let $V_k(\xi)$ be the Stiefel manifold bundle associated to $\xi$. 
It's a fiber bundle over $B$ with fiber $F=V_k(\mR^n)$. A section $s$ of 
$V_k(\xi)$ is a $k$-frame of $\xi$. Locally on a trivializing
neighborhood $U_\alpha$, $s$ can be identified as a map from $U_\alpha$ to
$V_k(\mR^n)$. Thus we can use obsruction theory to study the existence of
sections of $V_k(\xi)$. 

The obstructions $c^{j+1}(s)$ lie in $H^{j+1}(B;\mathcal{G})$, where 
$\mathcal{G}$ is a local coefficient system on $B$ with typical group 
$\pi_j(V_k(\mR^n))$. This is because although the ristriction of $V_k(\xi)$ 
on every $j$-cell $e^j_\alpha$ is trivial, the transition function may not 
be identity, from which induces a nontrivial isomorphism of homotopy groups. 
Moreover, any path $\gamma_{xy}$ from $x$ to $y$ can lift to a homotopy 
from $F_x$ to $F_y$, thus induces an isomorphism 
\begin{equation*}
    (\gamma_{xy})_*:\pi_j\big(F_x,x\big)
    \overset{\cong}{\to}\pi_j\big(F_y,y\big).
\end{equation*}
The theorems about obstruction classes can all generalize to cohomology with 
local coefficients without difficulty. 

Now since $V_k(\mR^n)$ is $(n-k-1)$-connected, all the section $s$ can be 
extended to $B^{(n-k)}$, all the sections on $B^{(n-k)}$ are homotopic with 
each other. The first obstruction comes when extend $s$ to 
$(n-k+1)$-skeleton, which we denote by 
\begin{equation*}
    \mathfrak{o}_{n-k+1}:=c^{n-k+1}(s)\in 
    H^{n-k+1}\Big(B;\big\{\pi_{n-k}\big(V_k(\mR^n)\big)\big\}\Big)
\end{equation*}
Since all sections on $B^{(n-k)}$ are homotopic with each other, this 
obstruction class does not depend on the choice of $s$, i.e, 
$\mathfrak{o}_{n-k+1}$ it is well-defined. 

Now we construct a series of characteristic classes from obsruction classes. 
\begin{itemize}
    \item If $n-k=2j-1$ is odd and $1<k<n$, then 
    $\pi_{2j-1}(V_{n-2j+1}(\mR^n)) = \mZ_2$. In this case, 
    \begin{align*}
        H^{2j}\Big(B;\big\{\pi_{2j-1}\big(V_{n-2j+1}(\mR^n)\big)\big\}\Big)
        &\cong H^{2j}(B;\mZ_2) \\ 
        \mathfrak{o}_{2j} &\mapsto w_{2j}(\xi).
    \end{align*}
    \item If $n-k=2j$ is even and $1<k<n$, then 
    $\pi_{2j}(V_{n-2j}(\mR^n)) = \mZ$. After applying the mod 2 reduction, 
    we have 
    \begin{align*}
        H^{2j+1}\Big(B;\big\{\pi_{2j}\big(V_{n-2j}(\mR^n)\big)\big\}\Big)
        &\xrightarrow{\text{mod } 2} H^{2j+1}(B;\mZ_2) \\ 
        \mathfrak{o}_{2j+1} &\mapsto w_{2j+1}(\xi).
    \end{align*}
    \item If $k=n$, then $V_n(\mR^n) = O(n)$, but $\pi_0(O(n)) = \mZ_2$ is 
    not a group. We turn to use $\tilde{H}^0(O(n);\mZ) = \mZ$ instead. Then 
    $\tilde{H}^0(O(n);\mZ)$ is a local coefficient system on $B$ with typical 
    group $\mZ$ and action induced by the determinant representation 
    $\det:\pi_1(B)\to O(n)\to \{\pm 1\}$. After applying the mod 2 reduction, 
    we have
    \begin{align*}
        H^1\Big(B;\big\{\tilde{H}^0(O(n))\big\}\Big)
        &\xrightarrow{\text{mod } 2} H^1(B;\mZ_2) \\ 
        \mathfrak{o}_1 &\mapsto w_1(\xi).
    \end{align*}
    \item If $k=1$, then $\pi_{n-1}(V_1(\mR^n)) = \pi_{n-1}(S^{n-1}) = \mZ$.
    After applying the mod 2 reduction, we have
    \begin{align*}
        H^n\Big(B;\big\{\pi_{n-1}(V_1(\mR^n))\big\}\Big)
        &\xrightarrow{\text{mod } 2} H^n(B;\mZ_2) \\ 
        \mathfrak{o}_n &\mapsto w_n(\xi).
    \end{align*}
\end{itemize}

\begin{center}
\begin{table}[htbp]
\begin{tabular}{|c|cc|}
\hline
Cases              & \multicolumn{1}{c|}{Obstruction classes}       & Characteristic classes       \\ \hline
$n-k=2j-1,\;1<k<n$ & \multicolumn{2}{c|}{$\mathfrak{o}_{2j}=w_{2j}$}                               \\ \hline
$n-k=2j,\;1<k<n$   & \multicolumn{2}{c|}{$\mathfrak{o}_{2j+1}\xrightarrow{\text{mod } 2}w_{2j+1}$} \\ \hline
$k=n$              & \multicolumn{2}{c|}{$\mathfrak{o}_{1}\xrightarrow{\text{mod } 2}w_{1}$}       \\ \hline
$k=1$              & \multicolumn{2}{c|}{$\mathfrak{o}_{n}\xrightarrow{\text{mod } 2}w_{n}$}       \\ \hline
\end{tabular}
\end{table}
\end{center}

\subsubsection*{Properties of our construction}
\noindent\textbf{Stiefel-Whitney Classes}
\begin{theorem}
    $w_i(\xi)\in H^i(B;\mZ_2)$ are the Stiefel-Whitney classes of $\xi$.
\end{theorem}
\begin{proof}
    omitted.
\end{proof}

\noindent\textbf{$\mathfrak{o}_n$ and the Euler Class}
If $\xi$ is orientable, then the structure group of $\xi$ can be 
reduced to $SO(n)$. Thus the action of $\pi_1(B)$ on 
$\pi_{n-k}(V_k(\mR^n))$ is trivial for $1\leq k<n$. So the local 
coefficient system is trivial. In this case, 
\begin{equation*}
    H^n\Big(B;\big\{\pi_{n-1}(V_1(\mR^n))\big\}\Big) 
    \cong H^n(B;\mZ).
\end{equation*}
\begin{theorem}
    If $\xi$ is orientable, then $\mathfrak{o}_n\in H^n(B;\mZ)$ 
    is the Euler class $e(\xi)$.
\end{theorem}
\begin{proof}
    omitted.
\end{proof}

\noindent\textbf{$w_1$ and Orientability}
The structure group of $\xi$ can be reduced to $SO(n)$ if and only if
the transition functions all have determinant 1, i.e., the action of 
$\pi_1(B)$ on $\tilde{H}^0(O(n);\mZ)$ is trivial. This is equivalent to
the triviality of $\mathfrak{o}_1$. Thus
\begin{proposition}
    $\mathfrak{o}_1 = 0$ if and only if $\xi$ is orientable.
\end{proposition}
\begin{theorem}
    $\mathfrak{o}_1$ has order 2, i.e., $2\mathfrak{o}_1 = 0$. Thus 
    \begin{equation*}
        w_1(\xi) = 0 \text{ if and only if } \mathfrak{o}_1 = 0.
    \end{equation*}
\end{theorem}
\begin{proof}
    The connected conponents of the space $V_n(\mR^n)$ are identified with 
    orientation of the fiber, and hence the elements of 
    $\tilde{H}^0(V_n(\mR^n);\mZ)\cong\mZ$ has the form 
    $a\mathrm{Or}_1-a\mathrm{Or}_2$; to each orientation we assign its 
    Coefficient.
    
    ...
\end{proof}

\noindent\textbf{Is $w_i$ the true obsruction?}
The fibration 
\begin{equation*}
    S^{2j-1}\to V_{n-2j+1}(\mR^n)\to V_{n-2j}(\mR^n).
\end{equation*}
induces an exact sequence of homotopy groups
\begin{equation*}
    \pi_{2j}(V_{n-2j}(\mR^n)) \xrightarrow{\partial} 
    \pi_{2j-1}(S^{2j-1})\to \pi_{2j-1} (V_{n-2j+1}(\mR^n)) \to 0.
\end{equation*}
It is actually 
\begin{equation*}
    \mZ \xrightarrow{\p} \mZ \xrightarrow{i_*} \mZ_2 \to 0.
\end{equation*}
Therefore $\mathrm{Im}\,\p=\Ker i_* = 2\mZ$ and $\Ker\p = 0$. So $\p$ is
multiplication by 2. So the exact sequence is a short exact sequence
\begin{equation*}
    0\to \mZ \xrightarrow{\times 2} \mZ \xrightarrow{i_*} \mZ_2 \to 0.
\end{equation*}
This short exact sequence induces a long exact sequence of cohomology groups
\begin{equation*}
    \cdots \to H^{2j}(B;\mZ) 
    \xrightarrow{i_*} H^{2j}(B;\mZ_2) \xrightarrow{\beta} H^{2j+1}(B;\mZ) 
    \xrightarrow{\times 2} H^{2j+1}(B;\mZ) \to \cdots
\end{equation*}
where $\beta$ is the Bockstein homomorphism.
\begin{theorem}
    $\mathfrak{o}_{2j+1} = \beta(\mathfrak{o}_{2j})$.
\end{theorem}
\begin{proof}
    omitted.
\end{proof}
\begin{corollary}
    $\mathfrak{o}_{2j+1}$ has order 2, i.e., 
    $2\mathfrak{o}_{2j+1} = 0$. Thus 
    \begin{equation*}
        w_{2j+1}(\xi) = 0 \text{ if and only if } 
        \mathfrak{o}_{2j+1} = 0.
    \end{equation*}
\end{corollary}
In summary, almost all the Stiefel-Whitney classes $w_i$ are the true 
obstructions, except when $i=n$ is even. If in addition $\xi$ is oriented, 
then $\mathfrak{o}_n$ is Euler class $e(\xi)$.